\documentclass[11pt,a4paper]{article}

\usepackage[margin=2.5cm]{geometry}
\usepackage{booktabs}
\usepackage{multirow}
\usepackage{amsmath}
\usepackage{graphicx}
\usepackage{caption}
\usepackage{hyperref}
\usepackage{parskip}

\hypersetup{
    colorlinks=true,
    linkcolor=black,
    citecolor=black,
    urlcolor=blue
}

\title{\textbf{Evaluation and Methodological Correction of Gated Feature Fusion in GTE-PPIS}}
\author{Pranav Nagaraji}
\date{August 2026}

\begin{document}

\maketitle

\begin{abstract}
This report documents the implementation of a Gated Feature Fusion Module (FFM) extending
the GTE-PPIS dual-branch architecture (Graph Transformer + Equivariant GNN) with
ESM-2 protein language model embeddings, the identification and correction of a
test-set threshold-selection data leak in the original evaluation pipeline, and the
resulting corrected benchmark results across three standard protein--protein interaction
site (PPIS) test datasets. All results reported here use a validation-locked threshold
protocol, ensuring that test labels are never used during decision boundary selection.
The work is a work-in-progress research extension and is not claimed to be
publication-ready.
\end{abstract}

\tableofcontents
\vspace{1em}

% ------------------------------------------------------------------
\section{Background and Motivation}
% ------------------------------------------------------------------

GTE-PPIS (Wang et al., 2025) is a dual-branch model combining a Graph Transformer
(GT) and an Equivariant Graph Neural Network (EGNN) for predicting protein--protein
interaction sites at the residue level. The published model uses 61-dimensional
handcrafted node features comprising DSSP secondary structure descriptors (14d),
PSSM evolutionary profiles (20d), HMM profiles (20d), and residue atom features (7d),
without any protein language model (PLM) integration.

This work extends GTE-PPIS by injecting ESM-2 (650M) per-residue embeddings through
a learned \emph{Gated Feature Fusion Module} (FFM) applied at the node level,
before graph message passing, to adaptively combine handcrafted evolutionary features
with PLM representations. This specific formulation --- a scalar sigmoid gate operating
on the classical vs.\ PLM embedding axis at the pre-GNN input stage --- has not
been applied in the established GTE-PPIS-family literature on the standard
Test\_60/Test\_315-28/UBtest\_31-6 benchmark family, and the closest related work
(DHEG; Li et al., 2026, \textit{Briefings in Bioinformatics}) explicitly identifies
effective fusion of PLM and geometric graph representations as an open problem.

% ------------------------------------------------------------------
\section{Architecture}
% ------------------------------------------------------------------

\subsection{Gated Feature Fusion Module}

For each residue $i$, let $\mathbf{c}_i \in \mathbb{R}^{40}$ denote the concatenated
PSSM and HMM features, and let $\mathbf{p}_i \in \mathbb{R}^{1280}$ denote the
ESM-2 embedding. The FFM projects both modalities into a shared space of dimension
$d = 128$ and computes a scalar gate:
%
\begin{align}
    \tilde{\mathbf{c}}_i &= \mathrm{LN}\bigl(W_c\,\mathbf{c}_i\bigr), \qquad
    \tilde{\mathbf{p}}_i = \mathrm{LN}\bigl(W_p\,\mathbf{p}_i\bigr), \\
    g_i &= \sigma\!\left(W_g\,[\tilde{\mathbf{c}}_i \,\|\, \tilde{\mathbf{p}}_i] + b_g\right), \\
    \mathbf{f}_i &= g_i \cdot \tilde{\mathbf{c}}_i + (1 - g_i) \cdot \tilde{\mathbf{p}}_i,
\end{align}
%
where $\sigma$ is the sigmoid function, $\|$ denotes concatenation, and
$\mathrm{LN}$ is Layer Normalization. The fused representation $\mathbf{f}_i \in \mathbb{R}^{128}$
is then concatenated with the DSSP (14d) and residue atom (7d) features to form
the 149-dimensional node input for the EGNN and GT branches.

\subsection{Stability Modifications}

Prior to the reported run, training was unstable with multiple folds collapsing.
Three targeted modifications were made:
\begin{enumerate}
    \item \textbf{EGNN residual connections enabled.} The 10-layer EGNN was
          previously instantiated with \texttt{residual=False}. Enabling
          \texttt{residual=True} adds skip connections inside each E\_GCL block,
          substantially improving gradient flow.
    \item \textbf{Focal Loss with per-fold class weighting.}
          The standard \texttt{CrossEntropyLoss} was replaced with Focal Loss
          (Lin et al., 2017):
          $\mathrm{FL}(p_t) = -\alpha_t(1-p_t)^\gamma \log p_t$,
          with $\gamma = 2.0$ and $\alpha = [1.0,\; w^+]$ where
          $w^+ = N^-/N^+$ is computed per fold from the training subset.
          The global ratio from Train\_335 is $w^+ = 5.4056$ (15.61\% positive
          residues across 66{,}208 total residues).
    \item \textbf{Degenerate threshold excluded.} The threshold search scan was
          changed from \texttt{range(0, 100)} to \texttt{range(1, 100)}, preventing
          the trivially valid threshold $\tau = 0.0$ from masking collapsed models.
\end{enumerate}

% ------------------------------------------------------------------
\section{Training Configuration}
% ------------------------------------------------------------------

Table~\ref{tab:config} summarises the training hyperparameters for the evaluated run.

\begin{table}[h]
\centering
\caption{Training configuration for run \texttt{fusion\_gated\_d128\_2026-07-30-12-18-16}.}
\label{tab:config}
\begin{tabular}{ll}
\toprule
\textbf{Parameter} & \textbf{Value} \\
\midrule
Fusion mode               & Gated FFM \\
Projection dimension ($d$) & 128 \\
Input feature dimension   & 149 (fused) \\
EGNN layers               & 10 \\
GT layers                 & 4 \\
Hidden dimension          & 256 \\
EGNN residual connections & Enabled \\
Coordinate update (tanh)  & Enabled \\
Attention heads           & 4 \\
Dropout                   & 0.1 \\
Loss function             & Focal Loss ($\gamma = 2.0$) \\
Class weight $[\alpha^-, \alpha^+]$ & $[1.0,\; \text{per-fold}]$ \\
Global $w^+$ (Train\_335) & 5.4056 \\
Optimiser                 & Adam \\
Learning rate             & $10^{-3}$ \\
LR scheduler              & ReduceLROnPlateau (factor 0.6, patience 5) \\
Gradient clipping         & $\ell_2$ norm $\le 1.0$ \\
Training epochs           & 50 \\
Batch size                & 1 \\
Cross-validation folds    & 5 \\
Random seed               & 2024 \\
Training dataset          & Train\_335 (334 proteins, after removing \texttt{2j3rA}) \\
\bottomrule
\end{tabular}
\end{table}

% ------------------------------------------------------------------
\section{Per-Fold Training and Validation Results}
% ------------------------------------------------------------------

Checkpoints were saved at the epoch with the highest \emph{validation} AUPRC
(the criterion used by the original GTE-PPIS codebase). Table~\ref{tab:valid}
shows the per-fold validation dataset statistics and the performance of the saved
checkpoint at the best epoch, as read from the training log.

\begin{table}[h]
\centering
\caption{Per-fold training split and validation checkpoint performance.
Metrics are evaluated on the held-out validation fold of Train\_335.}
\label{tab:valid}
\setlength{\tabcolsep}{5pt}
\begin{tabular}{ccccccc}
\toprule
\textbf{Fold} & \textbf{Train/Val} & \textbf{Neg/Pos} & \textbf{Best} & \textbf{Val}  & \textbf{Val}  & \textbf{Val} \\
              & \textbf{Proteins}  & \textbf{Weight}  & \textbf{Epoch}& \textbf{AUPRC}& \textbf{AUROC}& \textbf{$\tau^*$} \\
\midrule
1 & 267 / 67 & 5.5172 &  5 & 0.4893 & 0.7989 & 0.61 \\
2 & 267 / 67 & 5.3592 &  8 & 0.4053 & 0.7636 & 0.74 \\
3 & 267 / 67 & 5.5976 & 10 & 0.4867 & 0.8060 & 0.39 \\
4 & 267 / 67 & 5.2892 & 14 & 0.4530 & 0.8078 & 0.61 \\
5 & 268 / 66 & 5.2623 &  9 & 0.4940 & 0.8159 & 0.54 \\
\midrule
Mean & -- & -- & 9.2 & 0.4657 & 0.7984 & 0.578 \\
\bottomrule
\end{tabular}
\end{table}

$\tau^*$ denotes the classification threshold maximising F1 on the validation set
and is used to evaluate the corresponding test fold without any access to test labels.

% ------------------------------------------------------------------
\section{Evaluation Protocol}
% ------------------------------------------------------------------

\subsection{Original Protocol: Test-Set Threshold Leakage}
\label{sec:leak}

In the original \texttt{test.py}, after generating test predictions for a given
checkpoint, the evaluation called:
\begin{verbatim}
result = analysis(test_true, test_pred)   # no threshold supplied
\end{verbatim}
Inside \texttt{analysis()}, when no threshold is supplied, the code searches
$\tau \in \{0.01, 0.02, \ldots, 0.99\}$ and selects the value that maximises F1
\emph{on the test set} (using \texttt{test\_true}). This constitutes data leakage:
the ground-truth test labels are used to tune the classification boundary, which can
produce optimistic bias in Accuracy, Precision, Recall, F1, and MCC. AUROC and
AUPRC are unaffected because they are computed from continuous score rankings and
are threshold-independent.

\subsection{Corrected Protocol: Validation-Locked Threshold}
\label{sec:protocol}

The corrected evaluation follows the procedure below for each fold checkpoint:

\begin{enumerate}
    \item Load \texttt{Fold}$k$\texttt{\_best\_model.pkl}.
    \item Reconstruct the fold-$k$ validation set by replaying the same
          \texttt{KFold(n\_splits=5, shuffle=True)} with \texttt{seed=2024}
          on Train\_335 (identical to training).
    \item Obtain validation predictions; search $\tau \in [0.01, 0.99]$
          to maximise \emph{validation} F1. Call this $\tau_k$.
    \item Lock $\tau_k$.
    \item Evaluate the model on the held-out test benchmark; apply $\tau_k$
          as the fixed decision boundary. Test labels are never consulted
          during threshold selection.
\end{enumerate}
For full models (\texttt{Full\_model\_*.pkl}), the mean of the five fold thresholds
($\bar{\tau} = 0.578$) is used.

\noindent\textbf{Validation split correctness.} The identity of the reconstruction
was verified by checking that the validation AUROC and AUPRC reported by the
re-applied checkpoint match the corresponding values in the training log to four
decimal places across all five folds (see Table~\ref{tab:valid}).

% ------------------------------------------------------------------
\section{Corrected Test Benchmark Results}
% ------------------------------------------------------------------

Results in Tables~\ref{tab:t60}--\ref{tab:ub} were obtained using the corrected
validation-locked threshold protocol described in Section~\ref{sec:protocol}.
AUROC and AUPRC are threshold-independent and are identical under both protocols.
The 5-fold cross-validation average is computed over fold checkpoints 1--5 only;
full model results are reported separately.

\subsection{Test\_60}

\begin{table}[h]
\centering
\caption{Test\_60 benchmark results (13{,}144 residues).
$\tau^*$: validation-locked threshold. CV~Avg: mean over folds~1--5.}
\label{tab:t60}
\setlength{\tabcolsep}{4.5pt}
\begin{tabular}{lcccccccc}
\toprule
\textbf{Model} & $\tau^*$ & \textbf{AUROC} & \textbf{AUPRC} & \textbf{MCC} & \textbf{F1} & \textbf{Prec.} & \textbf{Rec.} & \textbf{Acc.} \\
\midrule
Fold 1 & 0.61 & 0.8075 & 0.4698 & 0.3695 & 0.4697 & 0.4663 & 0.4733 & 0.8313 \\
Fold 2 & 0.74 & 0.7894 & 0.4514 & 0.3468 & 0.4611 & 0.3598 & 0.6419 & 0.7632 \\
Fold 3 & 0.39 & 0.8079 & 0.4885 & 0.3822 & 0.4892 & 0.3953 & 0.6414 & 0.7885 \\
Fold 4 & 0.61 & 0.8095 & 0.4876 & 0.3913 & 0.4901 & 0.4759 & 0.5051 & 0.8341 \\
Fold 5 & 0.54 & 0.7906 & 0.4653 & 0.3541 & 0.4632 & 0.4264 & 0.5070 & 0.8145 \\
\midrule
\textbf{CV Avg} & 0.578 & \textbf{0.8011} & \textbf{0.4725} & \textbf{0.3688} & \textbf{0.4747} & 0.4247 & 0.5537 & 0.8063 \\
\midrule
Full (ep.~45) & 0.578 & 0.8075 & 0.4698 & 0.3695 & 0.4697 & 0.4663 & 0.4733 & 0.8313 \\
Full (ep.~9)  & 0.578 & 0.8100 & 0.4310 & 0.3399 & 0.4319 & 0.3034 & 0.7495 & 0.7202 \\
\bottomrule
\end{tabular}
\end{table}

\subsection{Test\_315-28}

\begin{table}[h]
\centering
\caption{Test\_315-28 benchmark results (60{,}376 residues).}
\label{tab:t315}
\setlength{\tabcolsep}{4.5pt}
\begin{tabular}{lcccccccc}
\toprule
\textbf{Model} & $\tau^*$ & \textbf{AUROC} & \textbf{AUPRC} & \textbf{MCC} & \textbf{F1} & \textbf{Prec.} & \textbf{Rec.} & \textbf{Acc.} \\
\midrule
Fold 1 & 0.61 & 0.7895 & 0.4008 & 0.3190 & 0.4259 & 0.3697 & 0.5022 & 0.8079 \\
Fold 2 & 0.74 & 0.7745 & 0.3943 & 0.2999 & 0.4092 & 0.2993 & 0.6466 & 0.7350 \\
Fold 3 & 0.39 & 0.7993 & 0.4270 & 0.3373 & 0.4368 & 0.3217 & 0.6798 & 0.7513 \\
Fold 4 & 0.61 & 0.8002 & 0.4232 & 0.3497 & 0.4518 & 0.3772 & 0.5632 & 0.8061 \\
Fold 5 & 0.54 & 0.7824 & 0.4059 & 0.3252 & 0.4323 & 0.3540 & 0.5552 & 0.7931 \\
\midrule
\textbf{CV Avg} & 0.578 & \textbf{0.7892} & \textbf{0.4102} & \textbf{0.3262} & \textbf{0.4312} & 0.3444 & 0.5894 & 0.7787 \\
\midrule
Full (ep.~45) & 0.578 & 0.7865 & 0.3954 & 0.3246 & 0.4309 & 0.3337 & 0.6080 & 0.7722 \\
Full (ep.~9)  & 0.578 & 0.8100 & 0.4310 & 0.3399 & 0.4319 & 0.3033 & 0.7495 & 0.7202 \\
\bottomrule
\end{tabular}
\end{table}

\subsection{UBtest\_31-6}

\begin{table}[h]
\centering
\caption{UBtest\_31-6 benchmark results (5{,}917 residues).}
\label{tab:ub}
\setlength{\tabcolsep}{4.5pt}
\begin{tabular}{lcccccccc}
\toprule
\textbf{Model} & $\tau^*$ & \textbf{AUROC} & \textbf{AUPRC} & \textbf{MCC} & \textbf{F1} & \textbf{Prec.} & \textbf{Rec.} & \textbf{Acc.} \\
\midrule
Fold 1 & 0.61 & 0.8185 & 0.4123 & 0.3663 & 0.4408 & 0.4501 & 0.4318 & 0.8683 \\
Fold 2 & 0.74 & 0.8034 & 0.3907 & 0.3530 & 0.4353 & 0.3383 & 0.6104 & 0.8097 \\
Fold 3 & 0.39 & 0.8212 & 0.4336 & 0.3752 & 0.4554 & 0.3672 & 0.5992 & 0.8278 \\
Fold 4 & 0.61 & 0.7943 & 0.3780 & 0.3155 & 0.3980 & 0.3966 & 0.3994 & 0.8548 \\
Fold 5 & 0.54 & 0.8038 & 0.4068 & 0.3647 & 0.4362 & 0.4615 & 0.4135 & 0.8716 \\
\midrule
\textbf{CV Avg} & 0.578 & \textbf{0.8083} & \textbf{0.4043} & \textbf{0.3549} & \textbf{0.4331} & 0.4028 & 0.4909 & 0.8464 \\
\midrule
Full (ep.~45) & 0.578 & 0.7887 & 0.3695 & 0.3130 & 0.4044 & 0.3438 & 0.4909 & 0.8263 \\
Full (ep.~9)  & 0.578 & 0.8301 & 0.4343 & 0.3674 & 0.4430 & 0.3310 & 0.6695 & 0.7977 \\
\bottomrule
\end{tabular}
\end{table}

% ------------------------------------------------------------------
\section{Observations}
% ------------------------------------------------------------------

\begin{itemize}
    \item All five fold checkpoints converge stably (MCC 0.30--0.39 across all
          three benchmarks), confirming that the stability fixes (EGNN residual
          connections and Focal Loss) successfully eliminated the training collapse
          observed in the prior run.
    \item The spread in locked validation thresholds ($\tau^* \in [0.39, 0.74]$)
          reflects natural between-fold class distribution variation and highlights
          why a fixed default threshold of 0.5 is inappropriate for this task.
    \item Fold~3, which received the lowest threshold (0.39), consistently achieves
          the highest recall across all three benchmarks, while Fold~2 (highest
          threshold, 0.74) achieves the highest precision. This precision--recall
          trade-off is expected under different decision boundaries and does not
          indicate a quality deficiency in either fold's model weights.
    \item Full model epoch~9 consistently matches or exceeds fold-average AUROC
          across all three benchmarks but at the cost of lower precision, suggesting
          early-stopped training benefits from the full dataset.
    \item A comparison against the original GTE-PPIS baseline and other published
          PPIS methods is deferred to future work when a complete ablation
          (\texttt{none}, \texttt{concat}, \texttt{gated}, \texttt{cross\_attn}) is
          available under the corrected evaluation protocol.
\end{itemize}

% ------------------------------------------------------------------
\section{Reproducibility}
% ------------------------------------------------------------------

The evaluation results in this report can be reproduced by running the following
command from the project root directory, assuming the virtual environment and
checkpoint directory are present:

\begin{verbatim}
source venv/bin/activate && python test.py \
    --fusion_mode gated \
    --d_proj 128 \
    --model_dir Log/fusion_gated_d128_2026-07-30-12-18-16/model/
\end{verbatim}

No retraining is required. The validation splits are deterministically reconstructed
inside \texttt{test.py} using \texttt{SEED = 2024} and the same \texttt{KFold}
configuration as \texttt{train.py}.

% ------------------------------------------------------------------
\section{Limitations and Next Steps}
% ------------------------------------------------------------------

\begin{enumerate}
    \item \textbf{Incomplete ablation.} The results reported here cover only the
          \texttt{gated} fusion mode. The full 4-condition ablation
          (\texttt{none} baseline, \texttt{concat}, \texttt{gated},
          \texttt{cross\_attn}) is required before any contribution of the FFM gate
          over simpler fusion strategies can be asserted.

    \item \textbf{No comparison against published baselines.} Direct numerical
          comparison against GTE-PPIS (Wang et al., 2025), DHEG (Li et al., 2026),
          and other PPIS methods requires careful protocol alignment (identical
          train/test splits, thresholding approach) and is outside the scope of
          this preliminary report.

    \item \textbf{Cross-attention arm.} The \texttt{cross\_attn} mode is implemented
          but not yet trained. Under the current residue-wise architecture, the
          cross-attention degenerates to a residual PLM projection; extending it
          to true cross-sequence attention would require architectural changes.

    \item \textbf{Hyperparameter sensitivity.} The Focal Loss parameter $\gamma$
          and the projection dimension $d$ have been fixed at $2.0$ and $128$,
          respectively. A systematic search may further improve stability and
          performance.

    \item \textbf{Full model threshold.} The full model threshold is currently set
          to the mean of fold thresholds ($\bar{\tau} = 0.578$). Ideally, a
          dedicated hold-out validation set should be reserved from the full
          training data to derive this threshold independently.
\end{enumerate}

% ------------------------------------------------------------------
\section{Conclusion}
% ------------------------------------------------------------------

This report covers the following:

\begin{enumerate}
    \item \textbf{What was implemented.} A Gated Feature Fusion Module (FFM) was
          implemented and integrated into the GTE-PPIS dual-branch architecture,
          enabling learned residue-level fusion of ESM-2 (650M) embeddings with
          classical evolutionary features (PSSM, HMM) before graph message passing.
          A new \texttt{loss.py} module implementing Focal Loss and per-fold
          class-weight computation was introduced.

    \item \textbf{What was trained.} Run \texttt{fusion\_gated\_d128\_2026-07-30-12-18-16}
          was trained for 50~epochs with 5-fold cross-validation on Train\_335
          (334 proteins), producing stable checkpoints across all five folds after
          enabling EGNN residual connections and Focal Loss.

    \item \textbf{What was changed in \texttt{test.py}.} The threshold-search loop
          was redesigned to derive the classification threshold $\tau_k$ exclusively
          from the corresponding fold's validation predictions, locking it before any
          contact with test data. A \texttt{get\_validation\_dataframes()} function
          was added to deterministically reconstruct the same fold splits used during
          training.

    \item \textbf{Why the change was necessary.} The original evaluation searched
          thresholds on the test labels themselves, constituting direct data leakage
          into discrete metrics (Accuracy, Precision, Recall, F1, MCC). This
          introduced potential optimistic bias and is methodologically incorrect for
          reporting research results.

    \item \textbf{Corrected evaluation outcome.} Under the corrected protocol, the
          gated FFM model achieves 5-fold cross-validation averages of
          AUROC~=~0.801, AUPRC~=~0.473, MCC~=~0.369 on Test\_60;
          AUROC~=~0.789, AUPRC~=~0.410, MCC~=~0.326 on Test\_315-28; and
          AUROC~=~0.808, AUPRC~=~0.404, MCC~=~0.355 on UBtest\_31-6.
          All fold models are stable, with no threshold degeneracy.

    \item \textbf{What remains before claiming a contribution.} The complete
          4-condition ablation study must be executed under the corrected evaluation
          protocol, and the gated FFM must demonstrate a statistically meaningful
          and consistent improvement over the \texttt{none} and \texttt{concat}
          baselines before any research contribution can be claimed.
\end{enumerate}

\end{document}
